\documentclass[11pt]{article}

\usepackage{amsmath,amssymb}
\usepackage[margin=1in]{geometry}

\title{A Sum-of-Squares Comparison}
\author{}
\date{}

\begin{document}

\maketitle



Assume
\[
\sum_i a_i=\sum_i a'_i,\qquad
a_1\ge a_i,\qquad a'_1\ge a'_i,\qquad a'_1>a_1,
\]
and assume additionally that
\[
a'_i\le a_i \qquad \text{for all } i\ne 1.
\]

Define
\[
r:=\sum_i (a'_i)^2-\sum_i a_i^2.
\]

We show that \(r\ge 0\). For \(i>1\), define
\[
e_i:=a_i-a'_i\ge 0.
\]
Since the two vectors have the same total sum,
\[
a'_1-a_1=\sum_{i>1}e_i=:d>0.
\]
Thus \(a'_1=a_1+d\), and for \(i>1\), \(a'_i=a_i-e_i\). Therefore
\[
\begin{aligned}
r
&=\sum_i (a'_i)^2-\sum_i a_i^2\\
&=(a_1+d)^2-a_1^2
  +\sum_{i>1}\left((a_i-e_i)^2-a_i^2\right)\\
&=2a_1d+d^2+\sum_{i>1}\left(-2a_i e_i+e_i^2\right)\\
&=2\sum_{i>1}e_i(a_1-a_i)+d^2+\sum_{i>1}e_i^2.
\end{aligned}
\]
Since \(a_1\ge a_i\) and \(e_i\ge 0\), every term on the right-hand
side is nonnegative. Hence
\[
r\ge 0.
\]
In fact, since \(d>0\), we have \(d^2>0\), so \(r>0\).

\end{document}
